%% kapa: exam=1 round=27 year=2016 period=1 form=A subject=civil_law subject_ko=민법 seq=29 number=29
%% source: Q-Net 감정평가사 기출문제 (cst00309) artlSeq=5169402 file=extracted/1008/5169402/제27회 감정평가사 1차 문제/제27회 감정평가사 1교시문제(A형).pdf
%% conversion: kordoc 4.24.51 -> tools/gen1.py (auto); review=proofread-sonnet; flags=none
\begin{question}{29}
부합에 관한 설명으로 옳지 않은 것은? (다툼이 있으면 판례에 따름)
\begin{choices}
  \item 부동산에 동산이 부합한 경우, 동산의 가격이 부동산의 가격을 초과하더라도 부동산의 소유자가 부합한 동산의 소유권을 취득한다.
  \item 동산끼리 부합된 경우, 주종을 구별할 수 없는 때에는 각 동산의 소유자가 부합 당시의 가액의 비율로 합성물을 공유한다.
  \item 토지 위에 건물이 신축 완공된 경우에 건물은 토지에 부합하지 않는다.
  \item 권원이 없는 자가 토지소유자의 승낙 없이 그 토지 위에 나무를 심은 경우, 특별한 사정이 없는 한, 토지소유자에 대하여 그 나무의 소유권을 주장할 수 있다.
  \item 건물의 임차인이 권원에 기하여 증축한 부분이 독립성을 가지면 증축된 부분은 부합되지 않는다.
\end{choices}
\end{question}
\answer{4}
